\chapter{Чи «ідеальне \texorpdfstring{$\psi$}{ψ}» дає ідеальні скаляри?}\label{ch:g02}

Метрика конкурсу (розд.~\ref{ch:g01}, формула~\eqref{eq:g01-S}) спирається на тезу
дизайнерів: розрив між добрим $R^2_\psi$ і поганими похідними скалярами неможливий за
побудовою, бо «ідеальна карта потоку заробляє їх усі». Теза правильна --- у точці, де
похибка $\psi$ дорівнює нулю. Цей розділ про те, що відбувається \emph{поруч}. Відповідь
глибокого розбору~D1 (2026-09-22): відображення $\psi\mapsto$ скаляри не ліпшицеве в
$L^2$, і карта з $R^2_\psi=0{,}99993$ вже може давати положення магнітної осі не краще
за середнє по вибірці.

\section{Фізика: чому скаляри чутливі}\label{sec:g02-physics}

Сім скалярів Consistency --- це функціонали $\psi$ трьох різних типів (\cite{Manual},
розд.~3--4):
\begin{itemize}
\item $\Rax,\Zax$ --- положення \emph{критичної точки}, де $\grad\psi=0$ (O-точка, магнітна вісь);
\item $\kappa$, $\dup$, $\dlo$, об'єм --- форма \emph{лінії рівня, що проходить крізь
  сідло} (X-точку), тобто останньої замкненої поверхні (LCFS);
\item $\li$ --- об'ємне середнє $\Bpol^2\propto|\grad\psi|^2$ у межах LCFS.
\end{itemize}
Скорер міряє похибку $\psi$ у нормі $L^2$. Але положення лінії рівня $\psi=\psib$ при
збуренні $\delta\psi$ зсувається на $\delta n\approx\delta\psi/|\grad\psi|$, а біля X-точки
$|\grad\psi|\to0$: теорема про неявну функцію вироджується. Отже, мала $L^2$-похибка
\emph{не обмежує} зсув скаляра; обмеження з'являється лише в нормах, що контролюють
похідні $\psi$ (для положення критичної точки потрібна мала похибка $\grad\psi$ у $C^0$, тобто норма
типу $C^1$; у 2D її не гарантує навіть $H^2$, потрібне $H^s$ з $s>2$), а таких членів метрика не має. Кількісно це описує \emph{показник підсилення}
$\alpha$:
\begin{equation}\label{eq:g02-alpha}
\frac{|f(\psi+\delta\psi)-f(\psi)|}{\sigma_f}\;\approx\;C\,e^{\alpha},\qquad
e=\frac{\|\delta\psi\|}{\|\psi-\bar\psi\|},\qquad R^2_\psi=1-e^2 .
\end{equation}
Для ліпшицевого в $L^2$ відображення $\alpha=1$: похибка скаляра спадає разом із похибкою
$\psi$. Якщо $\alpha<1$, то при $e\to0$ відношення «похибка скаляра~/ похибка $\psi$»
росте як $e^{\alpha-1}\to\infty$ --- це і є \emph{підсилення}. У реєстрах показник
записано просто як «$\alpha$», а менше за одиницю значення --- як ознаку підсилення.

\begin{derivation}[звідки береться $\alpha\approx\tfrac12$ (оцінка путівника; перевірено 2026-09-30, E36)]
Біля максимуму $\psi\approx\psi_0-\tfrac12\lambda r^2$, $\lambda$ --- власне значення
гесіана. \emph{Гладке} збурення з градієнтом $g$ зсуває максимум на $\Delta r=g/\lambda$,
тобто лінійно за амплітудою: $\alpha=1$. \emph{Шорстке} збурення (незалежний шум у
кожному вузлі з амплітудою $\eta$) діє інакше: дискретний максимум може перескочити в
будь-який вузол, де детермінований спад $\tfrac12\lambda r^2$ менший за $\eta$, тож
$\Delta r\sim\sqrt{2\eta/\lambda}$ --- тобто $\alpha=\tfrac12$ і $\Delta r\propto1/\sqrt\lambda$.
Скорер шукає вісь саме так: локальний максимум по вузлах плюс параболічне уточнення
(\texttt{fusion\_scoring/o\_point.py}). Ця оцінка пояснює кластер 0,53--0,55 для осі й
закон $1/\sqrt\lambda$ з E14, але вона ж передбачає перехід до $\alpha\to1$, коли
$\sqrt{2\eta/\lambda}$ стає меншим за крок сітки. Саме така залежність від $\lambda$
спростувала «універсальне $\alpha=\tfrac12$» (R6, нижче). Перевіркою цієї картини мала
стати гіпотеза C1. Перевірка 2026-09-30 (п.~\ref{sec:g02-c1}) картину \emph{підтвердила} для
осі: перехід $1\to\tfrac12$ стоїть там, де $\eta\approx\eta^{*}=\tfrac12\lambda h^2$ (E36). А от
«три кластери» виявилися властивістю вибірки (R10).
\end{derivation}

\section{Експеримент без навчання}\label{sec:g02-method}

Ідея розбору D1 проста: мережа не потрібна. Беремо \emph{істинні} карти $\psi$ з
демо-розрядів DIII-D, додаємо збурення контрольованої $L^2$-норми, перераховуємо скаляри
\emph{тим самим кодом скорера} і міряємо, наскільки вони зсунулися. Нормування обрано так,
щоб збігатися зі скорером: $e$ з \eqref{eq:g02-alpha} рахується по всьому набору кадрів,
тож $R^2_\psi=1-e^2$ точно; похибку скаляра ділимо на його популяційне стандартне
відхилення $\sigma_f$. Оскільки $R^2_f=1-\mathrm{MSE}/\sigma_f^2$, типова похибка порядку
$1\,\sigma_f$ означає, що $R^2$ цього скаляра впав приблизно до нуля --- прогноз не кращий
за середнє.

Дані: 60 кадрів із трьох розрядів DIII-D (203702--203704), сітка $65\times65$, по три
зерна на кадр, тобто 180 вимірів на кожну точку. Збурення трьох сімейств.

\code{ourtry/03-deep-dives/D1-psi-to-scalars/conditioning.py}{105}{131}{(сімейства збурень)}

\begin{itemize}
\item \emph{Рядки 105--115, \texttt{pert\_band}:} шум, спектр якого обмежено кільцем
  хвильових чисел $k_{\mathrm{lo}}\le|k|<k_{\mathrm{hi}}$. Генеруємо білий шум, беремо
  двовимірне ШПФ, множимо на маску кільця, повертаємося назад і нормуємо до одиничної
  норми (\texttt{\_normalise}). Вісім кілець від $|k|=1$ до частоти Найквіста 32 дають
  спектральну криву $S(k)$.
\item \emph{Рядки 118--123, \texttt{build\_pca}:} головні компоненти всього набору кадрів
  через SVD центрованої матриці. Функція повертає середнє і \emph{хвіст} ---
  рядки $V^{\top}$ після перших \texttt{n\_keep} (за замовчуванням 20).
\item \emph{Рядки 126--131, \texttt{pert\_pca\_tail}:} випадковий напрям у відкинутому
  хвості. Це рівно та похибка, яку робить декодер, що зрізає розклад на 20 компонентах
  (як бейзлайн PCA+Ridge стартового набору). Білий шум --- \texttt{pert\_white}, рядки 101--102.
\end{itemize}

\code{ourtry/03-deep-dives/D1-psi-to-scalars/conditioning.py}{150}{171}{(нормування і еталонні скаляри)}

\begin{itemize}
\item \emph{Рядки 150--158:} сітка й маска з \texttt{d3d\_envelope.npz}; масштаб
  $\|\psi-\bar\psi\|$ рахується по всіх кадрах разом --- рівно знаменник $R^2_\psi$
  скорера. \texttt{per\_frame\_scale} перераховує його на один кадр, щоб збурення
  окремого кадру з відносною амплітудою $e$ давало ту саму $e$ в пулі.
\item \emph{Рядки 160--167:} скаляри істини (\texttt{scalars\_of}, рядки 134--137, викликає
  \texttt{extract\_lcfs} і \texttt{derive\_frame} скорера) і їхні $\sigma_f$. Кадри, де
  якийсь скаляр не обчислився, відкидаються (\texttt{ok}).
\item \emph{Рядки 169--171:} база PCA будується один раз на всіх кадрах.
\end{itemize}

\code{ourtry/03-deep-dives/D1-psi-to-scalars/conditioning.py}{173}{197}{(розгортка за амплітудою і спектром)}

\begin{itemize}
\item \emph{Рядки 173--181:} для кожної амплітуди, кадру й зерна будуємо збурення,
  масштабуємо до $e\cdot$\texttt{per\_frame\_scale}, перераховуємо скаляри й записуємо
  $|\Delta f|/\sigma_f$. Зерно генератора залежить від кадру, номера зерна й $e$ ---
  прогін відтворюваний до останньої цифри.
\item \emph{Рядки 182--189:} медіана по 180 вимірах (\texttt{nanmedian} --- кадри, де LCFS
  не знайшлася, випадають; їхню кількість пише поле \texttt{n}).
\item \emph{Рядки 191--197:} два амплітудні скани (білий шум і хвіст PCA) на
  $e={}$0,003; 0,01; 0,03; 0,1; 0,2; 0,3 і спектральний скан при
  $e=0{,}1$, тобто $R^2_\psi=0{,}99$.
\end{itemize}

\code{ourtry/03-deep-dives/D1-psi-to-scalars/conditioning.py}{199}{211}{(підгонка показника $\alpha$)}

\noindent \emph{Рядки 201--209:} для білого шуму підганяємо пряму в координатах
$\ln e$--$\ln(\text{похибка})$; нахил --- це $\alpha$, перетин $b$ --- логарифм сталої $C$
з~\eqref{eq:g02-alpha}. Звідси точка беззбитковості $e^{*}=\exp(-b/\alpha)$, де похибка
досягає $1\,\sigma$, і $R^2_\psi{}^{*}=1-e^{*2}$ (рахує \texttt{plots.py}, пише
\texttt{break\_even.md}).

\section{Що вийшло}\label{sec:g02-results}

Прогін 2026-09-22 (60 кадрів, 3 зерна). Ми перезапустили його 2026-09-29 з
\texttt{--out} у тимчасову теку: усі медіани й показники збіглися до останньої цифри; на
машині, де збирався путівник, прогін тривав близько 57~с (у \texttt{findings.md} --- «$\approx30$~с»).

\begin{center}\small
\begin{tabular}{@{}lcccc@{}}
\toprule
скаляр & $\alpha$ & $e^{*}$ ($1\,\sigma$) & $R^2_\psi{}^{*}$ & похибка при $R^2_\psi=0{,}99$ \\
\midrule
$\Rax$ & 0,535 & 0,0084 & 0,99993 & 3,76$\,\sigma$ = 3,4~см \\
$\Zax$ & 0,547 & 0,0100 & 0,99990 & 3,53$\,\sigma$ = 4,3~см \\
$\kappa$ & 0,417 & 0,1149 & 0,98681 & 0,94$\,\sigma$ \\
$\dup$ (\texttt{tri\_top}) & 0,347 & 0,0249 & 0,99938 & 1,62$\,\sigma$ \\
$\dlo$ (\texttt{tri\_bot}) & 0,420 & 0,0022 & $\approx1{,}00000$ & 4,95$\,\sigma$ \\
об'єм & 0,695 & 0,0408 & 0,99834 & 1,87$\,\sigma$ = 0,45~м$^3$ \\
$\li$ & 0,696 & 0,1545 & 0,97612 & 0,74$\,\sigma$ \\
\bottomrule
\end{tabular}
\end{center}
\noindent Джерело: \texttt{results.json} (показники) і \texttt{break\_even.md} (решта).

\begin{established}{E1}
Відображення $\psi\mapsto$ похідні скаляри не ліпшицеве в $L^2$: усі
$\alpha\in[0{,}35;\,0{,}70]$; $\Rax$ досягає похибки $1\,\sigma$ при $R^2_\psi=0{,}99993$;
при $R^2_\psi=0{,}99$ похибка осі 3,4~см \cite{RegEst}. \emph{Уточнення 2026-09-30:}
$\alpha<1$ виміряно для \emph{білого} (повузлового) шуму; для гладкого збурення
$\alpha\approx0{,}69\text{--}1{,}08$ (п.~\ref{sec:g02-c1}).
\end{established}
\begin{established}{E2}
Показники розпадаються на три кластери: 0,53--0,55 (положення екстремуму),
0,35--0,42 (контур крізь сідло), 0,70 (інтеграли) \cite{RegEst}. \emph{Обмеження
2026-09-30:} розбиття стабільне лише на цих трьох розрядах; на 28 розрядах DIII-D його немає
(п.~\ref{sec:g02-c1}, R10).
\end{established}
\begin{established}{E3}
Хвіст PCA шкідливіший за білий шум тієї ж енергії: $\Zax$ --- 64,4$\,\sigma$ проти
6,3$\,\sigma$ при $e=0{,}30$ \cite{RegEst}.
\end{established}

\begin{figure}[ht]
\centering
\includegraphics[width=\textwidth]{g02_amplification.jpg}
\caption{Медіанна похибка семи скалярів (у одиницях $\sigma_f$) проти відносної похибки
$\psi$; зліва --- білий шум, справа --- хвіст PCA. Пунктир --- $1\,\sigma$, де $R^2$ скаляра
падає приблизно до нуля; верхня вісь --- $R^2_\psi$. Власний розрахунок, код:
\texttt{Our try/03-deep-dives/D1-psi-to-scalars/\{conditioning,plots\}.py}.}
\label{fig:g02-amp}
\end{figure}

Для масштабу: Stokolesov та ін. \cite{Stokolesov2025} повідомляють середнє зміщення
0,04~м для LCFS DIII-D, відновленої лише зі струмів котушок. Наші 3,4--4,3~см при
$R^2_\psi=0{,}99$ --- той самий порядок, тобто числа фізично осмислені.

\paragraph{Що в цих числах варто читати обережно.} Реєстр подає їх чесно, але кілька
деталей видно лише у \texttt{results.json}:
\begin{itemize}
\item $\sigma_f$ пораховано лише по трьох розрядах. Для $\Rax$ це 8,9~мм --- третина
  ширини комірки сітки (26,6~мм), тож «$1\,\sigma$» тут --- субпіксельний зсув. Нормовані
  похибки через це виглядають більшими, ніж були б на повному корпусі. Але $\alpha$ --- нахил
  у логарифмічних координатах, і від $\sigma_f$ він не залежить: $\sigma_f$ зсуває лише
  перетин. Тому якісний висновок E1 стійкий, а абсолютні значення подано окремо.
\item Колонка «похибка при $R^2_\psi=0{,}99$» --- значення \emph{підгонки}, а не
  виміру: для $\Rax$ підгонка дає 3,76$\,\sigma$, а виміряна медіана при $e=0{,}1$ ---
  4,11$\,\sigma$. Для $\dlo$ точка $e^{*}=0{,}0022$ лежить \emph{нижче} найменшої виміряної
  амплітуди 0,003 (там медіана вже 1,19$\,\sigma$), тобто це екстраполяція.
\item E3 справджується при великих $e$, але не рівномірно: при $e=0{,}003$ хвіст PCA
  зсуває вісь \emph{менше} за білий шум ($\Rax$: 0,15 проти 0,48$\,\sigma$), при $e=0{,}1$ ---
  майже так само (4,4 проти 4,1$\,\sigma$), і лише далі різниця вибухає (при $e=0{,}3$:
  23,5 проти 6,1 для $\Rax$, 64,4 проти 6,3 для $\Zax$). Англійський підсумок реєстру
  «у 10 разів шкідливіший» --- це саме $e=0{,}3$. При $e=0{,}3$ хвіст PCA ще й ламає
  пошук LCFS у 7 зі 180 вимірів (\texttt{n}~= 173).
\end{itemize}
Чому хвіст PCA такий руйнівний при великих $e$, реєстр не пояснює. Правдоподібна
версія: компоненти хвоста --- це напрями, у яких кадри ансамблю \emph{справді} відрізняються
на дрібному масштабі, тобто вони зосереджені там, де $\psi$ змінюється від кадру до
кадру (біля осі й межі), а не розмазані по всій сітці, як білий шум. Це інтерпретація, не
вимір.

\paragraph{Спектральна крива і гіпотеза C2.} При фіксованій $\|\delta\psi\|$ ($e=0{,}1$)
різні скаляри бояться різних мод (рис.~\ref{fig:g02-spec}): об'єм і $\dup$ --- низьких
хвильових чисел $|k|=1\text{--}4$ (глобальна форма), $\Rax$ --- середніх ($|k|\approx6\text{--}9$),
$\Zax$ і $\li$ --- $|k|\approx9\text{--}13$ ($\li$ містить $|\grad\psi|^2$, тож підсилює
високі $k$); $\kappa$ майже нечутлива до $k$. Наслідок: пласка $L^2$-втрата на $\psi$ ---
не та цільова функція, і з кривої $S(k)$ можна вивести спектрально зважену. Це й було
гіпотезою C2; перевірка 2026-09-30 її \emph{спростувала} (п.~\ref{sec:g02-c2}).

\begin{figure}[ht]
\centering
\includegraphics[width=0.72\textwidth]{g02_spectral_sensitivity.jpg}
\caption{Спектральна чутливість $S(k)$: медіанна похибка скалярів при збуренні в кільці
хвильових чисел, $e=0{,}10$ ($R^2_\psi=0{,}99$). Власний розрахунок, код:
\texttt{Our try/03-deep-dives/D1-psi-to-scalars/\{conditioning,plots\}.py}.}
\label{fig:g02-spec}
\end{figure}

Що з цього випливло для проєкту: якщо $R^2_\psi$ і Consistency розчіплюються так рано,
ще один $L^2$-член метрики нічого не лікує. Потрібна \emph{диференціальна} величина ---
нев'язка рівняння Ґреда--Шафранова, норма другого порядку, яка бачить саме ті моди, що
їх пропускає $L^2$. Це обґрунтування осі новизни №1 і рішення D6 (розд.~\ref{ch:g04}).
Декомпозиція похибки DeepONet \cite{Lanthaler2022} (кодування + апроксимація +
реконструкція) тут отримує уточнення: член реконструкції б'є не рівномірно, а по
найчутливіших напрямках --- принаймні при великих $e$.

\section{E14, R6 і C3: тест відновлено}\label{sec:g02-noscript}

До 2026-09-30 цей підрозділ називався «Пункти без відтворюваного коду»: синтетичний тест
2026-09-23 не зберегли (знахідка N5). Під час перевірки C1 його відновлено
(\texttt{c1/e14\_curvature.py}, \cite{RegEst,RegLog}).

\begin{itemcard}{E14}{встановлено (виміряно); \textbf{відтворено 2026-09-30}}
Зміщення осі під шумом $\propto1/\sqrt\lambda$ від кривизни: $|\Delta x|\sqrt\lambda$
стала в межах $\pm7\,\%$ на 16-кратному діапазоні $\lambda$. Відновлений тест дає для
$R$ при $\eta=10^{-2}$ і $\lambda=1\ldots16$ значення 0,0424--0,0474, тобто $\pm5{,}6\,\%$.
Умова --- режим стрибків, $\eta\gg\eta^{*}$. Для $Z$ за того самого $\eta$ закон не
виконується ($\pm19\,\%$): крок сітки по $Z$ --- 5~см, і для великих $\lambda$ режим стрибків
ще не настав.
\end{itemcard}
\begin{itemcard}{R6}{спростовано; \textbf{механізм відтворено 2026-09-30}}
«$\alpha=\tfrac12$ --- універсальний закон для положення критичної точки». Синтетика
2026-09-23 дала 0,687 при $\lambda=1$ і 0,841 при $\lambda=4$. Відновлений тест у
фіксованому вікні $\eta\in[10^{-5};10^{-2}]$ дає 0,746 / 0,819 / 0,884 / 0,938 / 0,976
для $\lambda=1,2,4,8,16$: $\alpha$ росте з $\lambda$, бо вікно захоплює дедалі більше
лінійного режиму. Точні числа 2026-09-23 не відтворюються --- вікно того тесту невідоме.
\end{itemcard}
\begin{conjecture}{C3}
Число обумовленості з E14 ($1/\sqrt\lambda$ в O-точці) корисне для стратифікації тесту.
\emph{Спростує:} розподіл $\lambda$ по корпусу виявиться вузьким --- стратифікувати нема за
чим. \emph{Стан 2026-09-30:} на 28 розрядах DIII-D $\lambda\in[0{,}13;\,3{,}6]$ (міжквартильний
розмах 0,73--1,58) --- розподіл широкий, тож головний спростовувач не спрацював. Але $\alpha$
осі між терцилями $\lambda$ змінюється слабко й непослідовно (0,47 / 0,54 / 0,54). Вердикту
ще немає \cite{RegConj}.
\end{conjecture}
\noindent Внутрішня напруга, про яку тут писало попереднє видання путівника («E14 каже
$1/\sqrt\lambda$ стала, а R6 --- що показник залежить від $\lambda$»), розв'язалася: обидва
твердження --- наслідки одного механізму E36, і тепер це перевіряється кодом.

\section{Перевірка C1 (2026-09-30): кластери --- властивість вибірки}\label{sec:g02-c1}

\paragraph{Навіщо.} Реєстр називав C1 найбазовішою гіпотезою розбору: якщо три кластери E2
справді відповідають трьом механізмам (екстремум, контур крізь сідло, інтеграл), D1 ---
теорія; якщо ні --- каталог чисел. \emph{Спростує:} синтетика з контрольованою геометрією дає
інші кластери, або кластери зникають при зміні вибірки розрядів \cite{RegConj}.

\paragraph{Як.} Спершу --- \emph{попередня реєстрація}: до будь-яких нових обчислень ми
записали в \texttt{c1/THEORY.md}, що саме рахує скорер для кожного скаляра, яке $\alpha$
очікуємо, що таке «кластер» (одновимірна k-means, силует $>0{,}5$) і за яких чисел C1
вважається підтвердженою чи спростованою. Туди ж записано конкурентну гіпотезу H-alt:
«положення екстремуму ($\Rax,\Zax$, а також $\dup,\dlo$ --- це $R$ у найвищій і найнижчій
точках контуру) дає $\approx\tfrac12$, а значення й інтеграли --- $\gtrsim0{,}8$». Далі чотири
кроки: відтворення D1; локальні моделі (параболоїд, «крапля» з X-точкою); 12 однонульових
рівноваг Серфона--Фрайдберга \cite{Cerfon2010} через увесь скорер; реальні дані --- 3 демо-розряди,
25 розрядів DIII-D з корпусу \cite{FEC} і 3 розряди MAST, з бутстрепом по \emph{розрядах}.

\code{ourtry/03-deep-dives/D1-psi-to-scalars/c1/alpha_tools.py}{149}{169}{(бутстреп по розрядах)}

\begin{itemize}
\item \emph{Рядки 152--153:} одиниця перевибірки --- розряд, а не кадр. Кадри одного розряду
  сильно корельовані, тож перевибірка кадрів дала б оманливо вузькі інтервали.
\item \emph{Рядки 158--161:} розряди витягаємо з поверненням, збираємо всі їхні кадри й
  перераховуємо медіанні криві та $\alpha$ тим самим \texttt{analyse}.
\item \emph{Рядки 162--169:} для кожної перевибірки записуємо $\alpha$ і розбиття семи
  скалярів на групи (найкраще за силуетом і примусово на три). Частка перевибірок, у
  яких розбиття збігається з C1, --- головне число вердикту.
\end{itemize}

\paragraph{Що вийшло.} Білий шум, $\alpha$ --- пряма через $e=0{,}003\ldots0{,}3$, як у D1;
частка 1000 перевибірок із розбиттям C1:
\begin{center}\small
\begin{tabular}{@{}lcc@{}}
\toprule
вибірка & відмінність від E2 & розбиття C1 \\
\midrule
DIII-D, 3 демо-розряди (як у D1) & --- & 96,7\,\% \\
DIII-D, 28 розрядів & $\li$: 0,70 $\to$ 0,47 (в інтервалі осі) & \textbf{0\,\%} \\
MAST, 3 розряди & $\Zax$: 0,55 $\to$ 0,71 & 3,3\,\% \\
Серфон--Фрайдберг, 12 форм & $\li$ 0,92, $\kappa$ 0,57 & 0,3\,\% \\
\bottomrule
\end{tabular}
\end{center}
Для гладкого збурення ($|k|<4$) усі $\alpha\approx0{,}69\text{--}1{,}08$, і кластерів немає ніде.
H-alt теж не пройшла (0--1,2\,\%).

\begin{figure}[ht]
\centering
\includegraphics[width=\textwidth]{g02_c1_alpha_samples.jpg}
\caption{Показники $\alpha$ (білий шум, оцінювач D1) для чотирьох вибірок із 95\,\% бутстреп-інтервалами;
сірі смуги --- три кластери E2. Власний розрахунок, код:
\texttt{Our try/03-deep-dives/D1-psi-to-scalars/c1/\{real\_sweep,cf\_xpoint,summarise\}.py}.}
\label{fig:g02-c1-samples}
\end{figure}

\begin{refuted}{R10 (була гіпотеза C1)}
«Три кластери $\alpha$ відповідають трьом механізмам». Розбиття стабільне лише на трьох
демо-розрядах, на яких його й побачили; на 28 розрядах DIII-D, на MAST і на синтетиці його
немає (0--3,3\,\%). Спростовано за критеріями, записаними до запуску \cite{RegEst}.
\end{refuted}

\paragraph{Що встановлено натомість.} Механізм для осі з рамки «Виведення» (п.~\ref{sec:g02-physics})
перевірено на параболоїді з кривиною $\lambda=1\ldots16$ на сітці DIII-D. Вісь шукає сама функція
скорера \texttt{find\_o\_point}.

\code{ourtry/03-deep-dives/D1-psi-to-scalars/c1/local_models.py}{50}{68}{(тест P1 / E14)}

\begin{itemize}
\item \emph{Рядки 57--61:} центр параболоїда випадково зсунуто в межах півкроку сітки --- інакше
  максимум завжди сидів би точно у вузлі, і параболічне уточнення нічого б не показало.
\item \emph{Рядки 62--64:} той самий генератор дає або білий шум (незалежний у кожному
  вузлі), або гладкий ($|k|<4$, нормований на ту саму амплітуду у вузлі); вісь
  шукає функція скорера до і після збурення.
\item \emph{Рядки 65--68:} медіана зсуву по 300 спробах для кожного $\eta$. Наступний
  рядок (поза фрагментом) рахує з цих медіан локальні нахили між сусідніми $\eta$.
\end{itemize}

\begin{figure}[ht]
\centering
\includegraphics[width=0.72\textwidth]{g02_c1_m1_crossover.jpg}
\caption{Локальний нахил медіанного $|\Delta\Rax|$ проти $\eta/\eta^{*}$, $\eta^{*}=\lambda h_R^2/2$:
суцільні --- білий шум, пунктир --- гладке збурення; кольори --- $\lambda=1\ldots16$. Власний
розрахунок, код: \texttt{c1/local\_models.py}, \texttt{c1/summarise.py}.}
\label{fig:g02-c1-m1}
\end{figure}

\begin{established}{E36}
Для білого шуму нахил похибки осі дорівнює 1, поки $\eta<\eta^{*}=\lambda h^2/2$ (параболічне
уточнення лінійне), і $\approx0{,}4\text{--}0{,}5$, коли $\eta\gtrsim\eta^{*}$ (дискретний
максимум стрибає). Криві для різних $\lambda$ лягають на одну в координатах $\eta/\eta^{*}$.
Для гладкого збурення нахил 1. Звідси $\alpha\approx\tfrac12$ осі в D1 --- режим стрибків на
сітці, а R6 --- наслідок фіксованого вікна $\eta$ \cite{RegEst}.
\end{established}

\paragraph{Що лишилося відкритим.} (1)~$\li$ на 28 реальних розрядах має $\alpha=0{,}47$, а на
синтетиці $\approx0{,}92\text{--}0{,}97$: на реальних кадрах $\li$ чутливе до чогось, чого немає в
рівновагах Соловйова. (2)~$\dup,\dlo$ дають $0{,}34\text{--}0{,}42$, нижче передбаченої $\tfrac12$;
імовірна причина --- $R$ береться в \texttt{argmax}~$Z$ серед 512 точок контуру, тобто змінюється
стрибками. (3)~Хвіст PCA на 28 розрядах дає розбиття C1 у 100\,\% перевибірок, але на MAST і
синтетиці --- ні; базис PCA сам залежить від вибірки, тож це спостереження, а не підтримка C1.

\section{Перевірка C2 (2026-09-30): вага не лікує розчеплення}\label{sec:g02-c2}

\paragraph{Що стверджувалось.} C2: спектрально зважена втрата, виведена з кривої $S(k)$, покращує
реальні моделі. \emph{Спростує:} бейзлайн, навчений із нею і без неї, має однаковий $S'$
\cite{RegConj}. $S'$ ще не реалізовано (член Calibration потребує невизначеності), тож до запуску
мірою записали \textbf{Consistency} (головна) і композитний \textbf{S} офіційного скорера.

\paragraph{Як.} Як і для C1, спершу --- попередня реєстрація (\texttt{c2/THEORY.md}): правило
виведення ваги, три контролі й критерії. Потім:
\begin{itemize}
\item $S(k)$ переміряно на 28 розрядах. У робочій точці моделей крива майже монотонно спадає
  з~$k$: найшкідливіші \emph{низькі} хвильові числа, а не середні, як у D1 при $e=0{,}1$.
\item Робоча точка: пласка PCA+Ridge на валідаційних розрядах має $R^2_\psi=0{,}940$, тобто
  $e_{\mathrm{op}}=0{,}245$; береться $S(k)$ при $e=0{,}3$. Ваги кілець спадають від 0,91 до 0,49.
\item Чотири «руки»: пласка MSE; вага $S(k)$; вага $H^1$ ($1+(k/k_c)^2$ з тим самим діапазоном,
  але зростаюча); ті самі числа $S(k)$, переставлені випадково.
\item Моделі: PCA+Ridge (вага діє через зважену PCA) і UNet\_Lite зі стартового набору (вага
  прямо у втраті, 3 зерна). Розбиття --- як у \texttt{SCORE.md}: навчання на потокових розрядах
  0--35, валідація 36--39, тест 60--67.
\end{itemize}

\code{ourtry/03-deep-dives/D1-psi-to-scalars/c2/train.py}{116}{127}{(зважена втрата у просторі Фур'є)}

\begin{itemize}
\item \emph{Рядки 121--123:} \texttt{rfft2} зберігає лише половину спектра дійсного поля, тож
  стовпці 1--32 входять удвічі (їхні спряжені пари не записано). Ділення на $(N^2)^2$ за
  рівністю Парсеваля робить пласку вагу $w\equiv1$ точно рівною звичайній MSE.
\item \emph{Рядки 125--127:} втрата --- зважена сума $|\widehat{\delta\psi}(k)|^2$. Та сама
  функція працює для всіх чотирьох рук, змінюється лише масив \texttt{w}: різниця між руками
  --- тільки у вазі, а не в коді.
\end{itemize}

\paragraph{Що вийшло.} Вісім тестових розрядів, UNet\_Lite --- середнє трьох зерен, 95\,\%
інтервали --- бутстреп по розрядах:
\begin{center}\small
\begin{tabular}{@{}llccc@{}}
\toprule
модель & вага & Consistency & $\Delta$ до flat [95\,\% ДІ] & $R^2_\psi$ \\
\midrule
PCA+Ridge & усі чотири & 0,2698 & $|\Delta|<10^{-4}$ & 0,090 \\
UNet\_Lite & flat & 0,179 & --- & 0,976 \\
UNet\_Lite & $S(k)$ & 0,168 & $-0{,}011$ [$-0{,}044$; $+0{,}012$] & 0,976 \\
UNet\_Lite & $H^1$ & 0,162 & $-0{,}017$ [$-0{,}070$; $+0{,}012$] & 0,975 \\
UNet\_Lite & shuffled & 0,177 & $-0{,}002$ [$-0{,}049$; $+0{,}019$] & 0,975 \\
\bottomrule
\end{tabular}
\end{center}

\begin{figure}[ht]
\centering
\includegraphics[width=\textwidth]{g02_c2_arms.jpg}
\caption{Consistency і композитний S для двох моделей і чотирьох вагових схем; дрібні точки ---
окремі зерна UNet\_Lite. Власний розрахунок, код:
\texttt{Our try/03-deep-dives/D1-psi-to-scalars/c2/\{train,evaluate,summarise\}.py}.}
\label{fig:g02-c2}
\end{figure}

PCA+Ridge на вагу не реагує: вага змінює лише базис PCA, а 50 компонент зважених і незважених
даних охоплюють майже той самий простір. У UNet\_Lite жодна вага не допомогла, а розкид між
зернами (Consistency від 0,114 до 0,231) більший за будь-яку різницю між руками
(рис.~\ref{fig:g02-c2}).

\begin{refuted}{R11 (була гіпотеза C2)}
«Спектрально зважена втрата, виведена з $S(k)$, покращує реальні моделі». Для обох моделей
інтервал $\Delta$Consistency містить нуль; $S(k)$ не краща ні за $H^1$, ні за перемішану вагу.
Спростовано за критеріями, записаними до навчання \cite{RegEst}.
\end{refuted}

\begin{established}{E37}
Розчеплення $R^2_\psi$ і Consistency видно й на навченій моделі. На тих самих восьми розрядах
UNet\_Lite має $R^2_\psi=0{,}976$, але Consistency лише 0,179 ($R^2$ для $\Zax$ дорівнює $-1{,}92$,
для $\kappa$ --- $-1{,}39$). У PCA+Ridge карта потоку майже марна ($R^2_\psi=0{,}090$), а
Consistency вища --- 0,270. Це теза E1, перевірена вже не штучними збуреннями \cite{RegEst}.
\end{established}

\paragraph{Що з цього випливає.} Перерозподіл $L^2$-похибки за хвильовими числами розчеплення не
лікує. Логічний наступний крок --- пряма мета на скаляри або диференціальний член (нев'язка ГШ),
тобто гіпотеза C4 (розд.~\ref{ch:g04}). \emph{Застереження:} одна архітектура, три зерна, вісім
тестових розрядів; діапазон ваг лише $1{,}85\times$ --- сильніша вага не перевірялась.

\section{Гіпотези і завдання}\label{sec:g02-tasks}

\begin{practicum}[R10, R11, C3 своїми руками]
\begin{enumerate}
\item \textbf{Відтворення.} У теці стартового набору:
  \texttt{.venv/bin/python "../../Our try/03-deep-dives/D1-psi-to-scalars/conditioning.py"
  --frames 60 --seeds 3 --out /tmp/d1} і порівняйте \texttt{/tmp/d1/results.json} з
  оригіналом (\texttt{--out} не дасть перезаписати файл проєкту).
\item \textbf{R10: відтворення.} Команди --- у \texttt{c1/README.md} (близько 6~хв на
  8 ядрах). Звірте частки розбиття C1 з таблицею п.~\ref{sec:g02-c1}.
\item \textbf{Загадка $\li$.} Чому на реальних кадрах $\alpha(\li)=0{,}47$, а на синтетиці
  $\approx0{,}95$? Порівняйте, де в кадрі лежить LCFS (сідло чи дотик до маски) і як часто
  вона перескакує під шумом. Сирі похибки є в \texttt{c1/real\_sweep\_DIII-D.npz}.
\item \textbf{Дискретність $\dup,\dlo$.} Замініть у копії \texttt{derive.\_shape\_from\_contour}
  \texttt{argmax} на параболічне уточнення вершини контуру. Чи підніметься $\alpha$ до $\tfrac12$?
\item \textbf{C3.} Кривизну $\lambda$ в O-точці для кожного кадру вже рахує
  \texttt{c1/real\_sweep.py} (\texttt{axis\_curvature}); розподіл широкий. Наступний крок ---
  стратифікувати за $\lambda$ не синтетичний шум, а залишки справжнього бейзлайну PCA+Ridge.
\item \textbf{R11: сильніша вага.} У \texttt{c2/train.py} візьміть $S(k)$ при $e=1{,}0$
  (діапазон ваг більший) або піднесіть ваги до степеня 3. Чи з'явиться ефект, більший за розкид
  між зернами? Скільки зерен потрібно, щоб його побачити?
\item \textbf{E37: пряма мета на скаляри.} Додайте до UNet\_Lite другу голову, що прогнозує сім
  скалярів Consistency (їхні еталонні значення рахує \texttt{local\_score.build\_reference}).
  Чи зросте Consistency, і чим доведеться заплатити в $R^2_\psi$?
\end{enumerate}
\end{practicum}
